% ==================== 章节7：周期尽头的质因数 ====================

\hirochapter{周期尽头的质因数：和补习组一起学习 Shor 算法}

\section*{序章：十五，三乘五，演示结束？}
\phantomsection
\addcontentsline{toc}{section}{序章：十五，三乘五，演示结束？}

初星学园开放日的前一天，仓本千奈抱着一卷海报纸走进空教室。纸卷比她的手臂还粗，展开后几乎占满整张课桌。海报中央印着一座层层收窄的金属圆筒，电缆与管线从顶端垂下来，蓝色指示灯在外壳上连成一圈。

花海佑芽围着海报转了半圈。

「好厉害！这是什么机器？」

「父亲大人替我找来的量子计算机照片。开放日用它做展板，一定很气派吧！」

「看起来好霸气！我也想试试！」

「正是如此！既然是能改变世界的机器，外形当然也要足够气派才行呢！」

「照片里最显眼的，其实是稀释制冷机。」

筱泽广趴在窗边的桌上，只抬起了眼睛。她面前放着开放日的演示说明书，标题是“用 Shor 算法分解整数 15”。

千奈的笑容僵了一下。

「制冷机？」千奈低头看了看照片，「那量子计算机在哪里呢？」

「量子芯片在里面。以超导量子比特为例，需要用氦同位素混合物工作的稀释制冷机，把芯片冷却到毫开尔文量级。不是把整台机器泡进液氮。」\cite{ch07ibmtech}

「那明天要把它搬来吗？」佑芽问。

「不。会先把展台压坏。」

「那小广，我们明天到底要演示什么？」佑芽把说明书翻到最后一页，又从最后一页翻回第一页，「要把十五分成质数，对吧？三乘五。完成！」

她用完成冲刺后的气势举起双手。

「恭喜花海同学！」千奈鼓起掌来，「本补习组尚未开始补习，就已经圆满完成任务了呢！」

「不。完全没有开始。」

广慢吞吞地坐直身体。仅仅是这个动作，她就像已经完成了一组高强度训练，安静地停了几秒。

「十五当然可以心算。选择十五，只是为了让我们能把算法的每一步都写出来。Shor 算法真正关心的是：如果一个数大到不能这样一眼看穿，怎样借助量子计算机找到它的质因数。」

「也就是说，要挑战很大的数字！」佑芽握紧拳头，「和挑战姐姐一样！」

「不太一样。不过，是很难的挑战。」

广的嘴角稍稍抬起。

「所以，很有趣。」

\section{把分解换成找周期}

广在白板上写下：

\begin{equation}
15 = 3 \times 5.
\end{equation}

「质数，是大于一、并且正因数只有一和它自身的整数。质因数分解，就是把一个合数写成质数乘积。比如十五分解成三乘五。」

「那十二就是二乘二乘三！」佑芽立刻回答。

「正解。」

「我也明白了！」千奈举起手，「一百就是十乘十！」

「十不是质数。」

「……看来补习现在正式开始了呢。」

广又在白板上写下两行方向相反的箭头：

\begin{align*}
3 \times 5 &\longrightarrow 15 && \text{很容易},\\
15 &\longrightarrow 3 \times 5 && \text{小规模容易，大整数可能很难}.
\end{align*}

「把两个大质数乘起来很快。可是只给出乘积，再想把那两个质数找回来，随着位数增加就会变得困难。RSA 的安全性，建立在经典计算机难以快速分解足够大的整数这一点上。」\cite{ch07shor}

「R……SA？」

「一种公钥密码方案。它可以参与通信加密和数字签名，但网络安全不是只靠这一件东西，更不能说 RSA 一旦受威胁，所有账户口令就会自动泄露。」

「那仓本家暂时不会破产？」

「不能从这件事直接推出破产。」

千奈松了口气，随后又觉得这句话似乎没有提供多少安慰。

「像是把面粉和水揉成面团很容易，想把面团恢复成原来的每一粒面粉就很困难！」佑芽说。

「方向不错。」

「那 Shor 算法就是把面团重新拆开？」

「它不直接拆。」广摇了摇头，「它换成另一个问题。」

她写下两个字：

\begin{center}
\textbf{找周期。}
\end{center}

广拿起粉笔，把整条路线压进一张图里：

\begin{figure}[htbp]
\centering
\begin{tikzpicture}[
  font=\scriptsize,
  >=Latex,
  node distance=0.48cm,
  shorstep/.style={draw=primary,thick,rounded corners=2pt,
    align=center,minimum width=9.6cm,minimum height=0.72cm},
  shorflow/.style={->,thick,draw=primary},
  shorretry/.style={->,thick,dashed,draw=accent,rounded corners=5pt}
]
\node[shorstep,fill=codebg] (input) {输入待分解的合数 $N$};
\node[shorstep,fill=green!14,below=of input] (choose)
  {经典预处理：随机选择 $1<a<N$，计算 $\gcd(a,N)$};
\node[shorstep,fill=orange!18,below=of choose] (order)
  {量子子程序：寻找 $a^x\bmod N$ 的周期（阶）$r$};
\node[shorstep,fill=blue!12,below=of order] (post)
  {经典后处理：验证 $r$，计算 $\gcd(a^{r/2}\!\pm1,N)$};
\node[shorstep,fill=yellow!20,below=of post] (result)
  {得到非平凡因数；若验证失败，则重新测量或更换 $a$};

\draw[shorflow] (input) -- (choose);
\draw[shorflow] (choose) --
  node[right,font=\tiny,text=primary] {$\gcd(a,N)=1$} (order);
\draw[shorflow] (order) -- (post);
\draw[shorflow] (post) -- (result);
\draw[shorretry] (result.west) -- ++(-0.72,0) |-
  node[pos=0.25,left,font=\tiny,text=accent] {重试} (choose.west);
\node[font=\tiny,text=primary,anchor=west] at ([xshift=0.14cm]choose.east)
  {$\gcd(a,N)>1$ 时直接得到因数};
\end{tikzpicture}
\caption{Shor 算法的经典--量子--经典流程。量子计算只承担找阶，前后仍需经典计算与验证。}
\label{fig:ch07-shor-flow}
\end{figure}

「这里的 $\gcd$ 是最大公约数。」广说，「也就是能同时整除两个整数的最大正整数。」

「为什么一开始就要算它？」

「假设我们要分解 $N$，随机选了一个满足 $1<a<N$ 的整数 $a$。如果
\begin{equation}
\gcd(a,N)>1,
\end{equation}
那我们已经撞上了 $N$ 的一个因数，根本不用启动量子电路。」

「抽签第一下就中大奖！」佑芽说。

「嗯。只是不能指望总有这种运气。如果 $\gcd(a,N)=1$，才进入真正的找周期。」

千奈仔细看着路线图。

「所以量子计算机只负责中间一段？」

「没错。Shor 算法并不是把所有步骤都丢进量子计算机。前处理和最后取最大公约数，都是经典计算。量子部分最重要的任务，只有找出周期 $r$。」

「来做十五的例子。令
\begin{equation}
N=15, \qquad a=2.
\end{equation}
因为 $\gcd(2,15)=1$，所以继续。现在依次计算 $2^x$ 除以十五的余数。」

广把笔递给佑芽。佑芽在白板上飞快地写了起来。

\begin{center}
\small
\begin{tabular}{ccc}
\toprule
指数 $x$ & $2^x$ & $2^x \bmod 15$ \\
\midrule
0 & 1   & 1 \\
1 & 2   & 2 \\
2 & 4   & 4 \\
3 & 8   & 8 \\
4 & 16  & 1 \\
5 & 32  & 2 \\
6 & 64  & 4 \\
7 & 128 & 8 \\
\bottomrule
\end{tabular}
\end{center}

\begin{figure}[htbp]
\centering
\begin{tikzpicture}[
  font=\small,
  >=Latex,
  state/.style={circle,draw=primary,thick,fill=codebg,
    minimum size=0.82cm,inner sep=1pt},
  cycleflow/.style={->,very thick,draw=accent,bend left=18}
]
\node[state] (one) at (0,1.55) {$1$};
\node[state] (two) at (2.15,0) {$2$};
\node[state] (four) at (0,-1.55) {$4$};
\node[state] (eight) at (-2.15,0) {$8$};
\draw[cycleflow] (one) to node[above right,font=\scriptsize] {$\times2$} (two);
\draw[cycleflow] (two) to node[below right,font=\scriptsize] {$\times2$} (four);
\draw[cycleflow] (four) to node[below left,font=\scriptsize] {$\times2$} (eight);
\draw[cycleflow] (eight) to node[above left,font=\scriptsize] {$\times2$} (one);
\node[align=center,font=\scriptsize\bfseries,text=primary] at (0,0)
  {$\bmod 15$\\周期 $r=4$};
\end{tikzpicture}
\caption{$2^x\bmod15$ 的余数环。每乘一次 $2$ 前进一步，四步后回到 $1$。}
\label{fig:ch07-mod15-cycle}
\end{figure}

「一、二、四、八，然后又回到一！」佑芽用笔圈住两组相同的余数，「每四步一圈。就像绕操场！」

「对。这个最小的正周期四，就是二模十五的阶。一般地，当 $a$ 与 $N$ 互质时，我们寻找最小的正整数 $r$，使得
\begin{equation}
a^r \equiv 1 \pmod N.
\end{equation}
这里的 $\equiv$ 表示同余。$a^r \equiv 1 \pmod N$ 的意思是，$a^r$ 除以 $N$ 的余数等于一。」

「刚才的例子就是
\begin{equation}
2^4=16\equiv 1 \pmod {15},
\end{equation}
所以 $r=4$！」佑芽说。

「完全正确。」

\section{周期怎样交出因数}

广把 $r=4$ 写到白板中央。

「找周期不是终点。接下来才是它和因数之间的连接。因为 $a^r\equiv1\pmod N$，所以 $N$ 能整除 $a^r-1$。如果 $r$ 是偶数，就能作平方差分解：」

\begin{equation}
a^r-1=\left(a^{r/2}\right)^2-1
=\left(a^{r/2}-1\right)\left(a^{r/2}+1\right).
\end{equation}

「也就是说，$N$ 整除右边两个数的乘积。只要 $a^{r/2}$ 模 $N$ 既不等于 $1$，也不等于 $-1$，这两个括号通常会分别藏着 $N$ 的非平凡因数。用最大公约数就能把它们取出来。」

代入 $N=15,a=2,r=4$：

\begin{equation}
a^{r/2}=2^2=4.
\end{equation}

于是：

\begin{align}
\gcd(4-1,15)&=\gcd(3,15)=3,\\
\gcd(4+1,15)&=\gcd(5,15)=5.
\end{align}

因此：

\begin{equation}
\boxed{15=3\times5}
\end{equation}

「回来了！」佑芽指着最开始的式子，「虽然绕了特别大一圈，最后还是三乘五！」

「而且十五的因数，从周期四里面出现了。」千奈看看白板，又看看自己画的登山路线，「像是走到山顶以后，发现宝箱一直藏在起点的地板下面呢。」

「Shor 算法的确很绕。」广看着满满一白板的公式，语气听上去十分满足，「所以，适合认真学习。」

\section{量子计算机怎样找周期}

佑芽抱起手臂。

「可是刚才的周期，也是我们一项一项算出来的。数字很大的话，照样要跑很多圈吧？」

「经典计算如果只是逐项试，周期可能很长。量子计算的优势，正是在这里提取周期结构。」

广在白板上画出上下两排线。

「可以把量子部分想成两个寄存器。」

\begin{itemize}[leftmargin=2em]
  \item \textbf{计数寄存器}有 $m$ 个量子比特，可以表示 $Q=2^m$ 个指数 $x$；
  \item \textbf{工作寄存器}至少要能编码 $0$ 到 $N-1$ 的余数，所需量子比特数通常与 $\lceil\log_2N\rceil$ 同阶。
\end{itemize}

「两个寄存器分工不同，大小也不必相等。先把工作寄存器初始化为 $\lvert1\rangle$。」

「第一步，对计数寄存器施加 Hadamard 门，把它准备成许多 $x$ 的均匀叠加。设寄存器能表示 $0$ 到 $Q-1$，理想化地写就是：」

\begin{equation}
\frac{1}{\sqrt{Q}}\sum_{x=0}^{Q-1}\lvert x\rangle.
\end{equation}

「所谓叠加，是不是说它一次同时跑了从零到 $Q-1$ 的所有圈？」佑芽问。

「可以暂时这样想，但不能把它理解成‘最后能一次读出所有答案’。量子态一经测量，只会给出一个经典结果。算法要做的，是在测量之前让不同分支发生干涉，把我们需要的结构集中到少数结果上。」

第二步，让两个寄存器共同经过\textbf{受控模幂运算}：

\begin{equation}
\lvert x\rangle\lvert 1\rangle
\longmapsto
\lvert x\rangle\lvert a^x \bmod N\rangle.
\end{equation}

于是，两个寄存器的联合状态变成

\begin{equation}
\frac{1}{\sqrt{Q}}\sum_{x=0}^{Q-1}
\lvert x\rangle\lvert a^x\bmod N\rangle.
\end{equation}

对于 $N=15,a=2$，工作寄存器里的数值呈现

\begin{equation}
1,2,4,8,1,2,4,8,\ldots
\end{equation}

的周期结构。

「第三步，在计数寄存器上施加\textbf{逆量子傅里叶变换}，也就是 inverse QFT。」

「听起来像必杀技！」佑芽说。

「从效果看，确实有一点像。」广说，「受控模幂、相位估计和逆 QFT 合在一起，让不匹配的振幅相互抵消，让与周期相符的振幅相互加强。量子计算机没有挨个试除质数，也不会把所有因数一次读出来；它把余数序列中的周期结构，变成测量分布里的峰。」\cite{ch07ibmtutorial}

千奈赶紧在海报上补了几座尖尖的山峰。

「如果 $Q$ 取二的幂，并且精度足够，测量计数寄存器会得到一个 $m$ 位比特串。把这个比特串当作二进制整数 $y$，就有较大概率满足
\begin{equation}
\frac{y}{Q}\approx\frac{s}{r},
\end{equation}
其中 $s$ 是从零到 $r-1$ 的某个整数。之后用经典的\textbf{连分数展开}，从 $y/Q$ 的有理数近似中恢复 $s/r$，它的分母就是周期 $r$ 的候选值。」

「为什么是‘猜’？」千奈问。

「因为一次测量未必给出足够的信息。比如测到 $s=0$，只得到零；或者 $s$ 与 $r$ 有公因数，约分后只能看到 $r$ 的一个因数。我们要把候选 $r$ 代回
\begin{equation}
a^r\equiv1\pmod N
\end{equation}
进行验证。不对就再测一次。」

为了让图像更具体，广取 $Q=256$。在 $r=4$ 的理想示例中，明显的测量位置对应：

\begin{center}
\small
\setlength{\tabcolsep}{4pt}
\begin{tabular}{c c L{3.0cm} L{3.8cm}}
\toprule
可能的 $y$ & $y/Q$ & 连分数结果 & 候选分母 \\
\midrule
0   & $0$   & $0/1$               & $1$，没有帮助 \\
64  & $1/4$ & $1/4$               & $4$ \\
128 & $1/2$ & $2/4$ 约分为 $1/2$ & $2$，只看到真周期的约数 \\
192 & $3/4$ & $3/4$               & $4$ \\
\bottomrule
\end{tabular}
\end{center}

\begin{figure}[htbp]
\centering
\begin{tikzpicture}[x=2.15cm,y=2.0cm,font=\scriptsize,>=Latex]
\draw[->,thick,draw=primary] (-0.45,0) -- (3.75,0)
  node[right] {测量整数 $y$};
\draw[->,thick,draw=primary] (-0.25,-0.02) -- (-0.25,1.28);
\node[rotate=90,text=primary] at (-0.50,0.63) {测量概率 $P(y)$};
\foreach \x/\value/\phase in {
  0/0/0,
  1/64/{1/4},
  2/128/{1/2},
  3/192/{3/4}}
{
  \draw[very thick,draw=accent] (\x,0) -- (\x,1.02);
  \fill[accent] (\x,1.02) circle (1.5pt);
  \ifnum\x=0
    \node[below=3pt,xshift=2pt] at (\x,0) {$\value$};
  \else
    \node[below=3pt] at (\x,0) {$\value$};
  \fi
  \node[above=2pt,text=primary] at (\x,1.02) {$y/Q=\phase$};
}
\draw[dashed,draw=black!30] (0,0.51) -- (3,0.51);
\end{tikzpicture}
\caption{$Q=256$、$r=4$ 时的理想测量峰。一次测量只返回一个 $y$；重复运行后，峰会集中在 $0,64,128,192$ 附近。}
\label{fig:ch07-phase-peaks}
\end{figure}

\FloatBarrier

「所以一次没有测到四，不代表算法错了。」广说，「概率算法允许失败。重复、验证，是算法本身的一部分。」

「和跳高一样！」佑芽拍了一下手，「第一次没过杆，调整以后再跳！」

「是。只是量子态不会因为热血而跳得更高。」

「一点点也不会吗？」

「不会。」

「这就是 Shor 算法？」千奈问。

「是框架。实现时最费资源的通常是可逆的受控模乘、模幂电路，以及把它们放进可靠的量子纠错体系。纸上的 $\lvert x\rangle\lvert1\rangle\mapsto\lvert x\rangle\lvert a^x\bmod N\rangle$ 只有一行，真正编译成量子门，会麻烦得多。」

「一行说明，很多训练。」佑芽点头，「我懂。就像训练计划上只写着‘跑步’两个字。」

「对。然后会跑很久。」

广说完，似乎仅仅想象了一下跑步就累了。她把额头轻轻靠在白板边缘。

\section{失败也在算法之内}

千奈举起手。

「广同学，既然最后可能失败，那它还是一个厉害的算法吗？」

「当然。随机算法和量子算法经常用成功概率描述结果。单次失败不可怕，只要成功概率有可靠下界，并且能快速验证答案，就可以通过重复把整体成功率提高。」

广列出常见的失败情形：

\begin{itemize}[leftmargin=2em]
  \item 量子测量得到的信息不足，无法恢复正确的 $r$；
  \item 得到的周期 $r$ 是奇数，不能使用平方差的中间步骤；
  \item $r$ 虽为偶数，但 $a^{r/2}\equiv-1\pmod N$，最大公约数只给出 $1$ 或 $N$；
  \item 连分数给出的是错误候选或真周期的约数，代回验证不通过。
\end{itemize}

「处理方法通常很朴素：重新测量，或者换一个 $a$。」

「不断失败，不断再来。」千奈握住画笔，「真是一种非常鼓舞人的算法呢！」

「我也这么觉得。」

广抬起头。那一瞬间，她脸上的笑意比讲出正确答案时明显得多。

「不过广同学喜欢失败的理由，似乎和算法不太一样呢……」

\section{今天还不能拆掉 RSA}

为了防止佑芽第二天在台上宣布“量子计算机能瞬间破解所有密码”，广在演示稿末尾画了一个很大的方框。

\begin{center}
\setlength{\fboxsep}{8pt}%
\colorbox{codebg}{%
\parbox{0.88\linewidth}{%
Shor 算法展示的是\textbf{理论上的高效}，不等于今天的量子计算机已经能轻易分解现实中的超大整数。}}
\end{center}

「更准确地说，Shor 算法的运行时间随输入位数，也就是 $\log N$，按多项式增长。相较于目前已知的经典通用整数分解算法，这是根本性的复杂度优势。但它仍然需要大量量子门、足够多的纠错量子比特，以及可靠的容错运行。」\cite{ch07shor,ch07ibmtutorial}

「那现在能分解十五吗？」佑芽问。

「能做小规模、经过编译或简化的演示。十五也确实是历史上最著名的实验示例之一。但是十五本来就能心算，这种实验的意义是验证控制、相位估计和量子电路，而不是证明今天已经能攻击现代 RSA 密钥。」\cite{ch07ibmtutorial}

「而且它也不是‘所有密码的克星’。」广继续写道，「Shor 算法直接威胁的是依赖整数分解或离散对数困难性的公钥体系。对称加密、哈希、口令保存等问题，不能简单说成被同一种方式全部破解。」

千奈在大方框旁边画了一只举着盾牌的小狮子。

「也就是说，不能夸大其词，但也不能装作它无事发生。」

「很好。」

「我明天就这样讲！」

「千奈，这一句本来是我的台词吧？」佑芽问。

「那花海同学负责最有气势地宣布答案！」

「好！十五等于三乘五——！」

声音撞上空教室的墙壁，又完整地弹了回来。广身体一晃，伸手扶住白板。

「佑芽。音量……不属于量子加速的一部分。」

\section{开放日}

第二天，三人的展台前很快围起了一圈观众。

千奈负责图画。她把模幂序列画成一条环形跑道：一、二、四、八四名选手不断绕圈；又把逆量子傅里叶变换画成聚光灯，让分散的波纹在四座山峰上相互加强。尽管气派的量子计算机并没有被搬到会场来，但不妨碍小孩子们看得目不转睛。

佑芽负责示范。她沿着地上贴出的数字格奔跑，每到第四格就举起写着“余数一”的牌子，然后用全场都能听见的声音宣布：

「找到周期四！接下来用最大公约数，把三和五从十五里面拉出来！」

广负责讲解。她先说明质因数分解，再把问题转为找阶；从叠加态讲到受控模幂，从干涉讲到逆 QFT，最后停在连分数和经典后处理。她的声音不大，语速也慢，却总能在听众即将迷路之前，用一句话指出路线：

「不要问量子计算机‘因数是什么’。要问它‘余数以什么周期重复’。周期知道以后，因数会从平方差里出现。」

演示结束，掌声响起。

广扶着桌沿坐下，像是她才是刚才沿数字格跑了十圈的人。

「小广，没事吧？」

「没事。只是站了二十分钟，腿已经不太听话。」

「非常辛苦呢。」千奈把水杯递给她，「不过大家似乎都听懂了！」

「嗯。」

广接过水杯，看向海报上的量子计算机。千奈在稀释制冷机旁边添了一张柔软的椅子，还用箭头标着“讲解员休息处”。

「白板上的计算很简单。把它讲到让人明白，比算出答案难。」

「那广开心吗？」佑芽问。

「很开心。」

她轻轻笑了。

「下次，讲离散对数吧。」

「还有下次吗！？」

千奈与佑芽的声音重叠在一起，撞上展板后慢了一拍，又传回三人耳边，仿佛某个终于被测量出来的周期。

\section{数学小结}

以 $N=15,a=2$ 为例，Shor 算法的关键链条是：

\begin{enumerate}[leftmargin=2.4em]
  \item $\gcd(2,15)=1$，进入找阶步骤。
  \item $2^x\bmod15$ 的序列为 $1,2,4,8,1,\ldots$，故阶为 $r=4$。
  \item $r$ 为偶数，并且 $2^{r/2}=4\not\equiv-1\pmod{15}$。
  \item 计算
  \begin{align*}
    \gcd(2^{r/2}-1,15)&=3,\\
    \gcd(2^{r/2}+1,15)&=5.
  \end{align*}
  \item 得到 $15=3\times5$。
\end{enumerate}

量子加速发生在第 2 步的\textbf{周期/阶寻找}；其余主要是经典数论计算。

\begin{thebibliography}{9}

\bibitem{ch07ibmlearning}
IBM Quantum Learning. \emph{Shor's algorithm}.
\url{https://quantum.cloud.ibm.com/learning/en/courses/fundamentals-of-quantum-algorithms/phase-estimation-and-factoring/shor-algorithm}.

\bibitem{ch07ibmtutorial}
IBM Quantum Documentation. \emph{Shor's algorithm tutorial}.
\url{https://quantum.cloud.ibm.com/docs/en/tutorials/shors-algorithm}.

\bibitem{ch07ibmtech}
IBM Quantum Learning. \emph{IBM Quantum technology}.
\url{https://quantum.cloud.ibm.com/learning/en/courses/quantum-business-foundations/quantum-technology}.

\bibitem{ch07shor}
P. W. Shor. Polynomial-time algorithms for prime factorization and discrete logarithms on a quantum computer.
\emph{SIAM Review}, 41(2):303--332, 1999.
\url{https://arxiv.org/abs/quant-ph/9508027}.

\end{thebibliography}
